\section{Verificación de requisitos funcionales}

La Tabla \ref{tab:verificacion_rf} presenta la correspondencia entre los requisitos funcionales definidos en el diseño y los módulos que materializan cada uno en la implementación. Todos los requisitos de prioridad alta y media definidos en la especificación fueron implementados en el prototipo entregado.

\begin{table}[H]
\centering
\caption{Verificación del cumplimiento de los requisitos funcionales}
\label{tab:verificacion_rf}
\footnotesize
\setlength{\tabcolsep}{4pt}
\begin{tabular}{
  >{\centering\arraybackslash}p{1.2cm}
  >{\raggedright\arraybackslash}p{3.8cm}
  >{\raggedright\arraybackslash}p{6.8cm}
  >{\centering\arraybackslash}p{2.1cm}
}
\hline
\textbf{ID} & \textbf{Nombre} & \textbf{Módulos de implementación} & \textbf{Estado} \\
\hline
RF-01 & Generación del AST &
  \texttt{\_runner.rkt}, \texttt{\_stream-parser.rkt}, \texttt{racket.ts} &
  Implementado \\
RF-02 & Visualización del AST &
  \texttt{ASTViewer.tsx}, \texttt{PlaygroundLayout.tsx} &
  Implementado \\
RF-03 & Ejecución y evaluación &
  \texttt{app/api/run/route.ts} (\textit{endpoint} \texttt{/api/run}) &
  Implementado \\
RF-04 & Representación del ambiente &
  \texttt{\_tracking.rkt}, \texttt{EnvironmentPanel.tsx} &
  Implementado \\
RF-05 & Biblioteca de ejemplos &
  \texttt{examples/}, \texttt{load-examples.ts}, \texttt{WelcomeModal.tsx} &
  Implementado \\
RF-06 & Validación de sintaxis Racket &
  \texttt{bnf-lexer.ts}, \texttt{bnf-parser.ts}, \texttt{cleanStderr()} en \texttt{route.ts} &
  Implementado \\
RF-07 & Visualización del intérprete &
  \texttt{EditorPanel.tsx}, \texttt{CodeEditor.tsx} (líneas bloqueadas con Monaco Editor) &
  Implementado \\
RF-08 & Plantilla base y exportación &
  \texttt{INITIAL\_FILES} y \texttt{downloadZip()} en \texttt{playground-utils.ts} &
  Implementado \\
\hline
\end{tabular}
\vskip 4pt
\footnotesize\textit{Fuente: elaboración propia.}
\end{table}

\FloatBarrier
